  \item[2.52] (a) $+\SI{1.5}{\metre\per\second\squared}$ for \SI{2}{\second}, $0$ for \SI{6}{\second}, $\SI{-1.0}{\metre\per\second\squared}$ for \SI{3}{\second}: three horizontal steps. (b) Triangle $+$ rectangle $+$ triangle: $3.0 + 18 + 4.5 = \SI{25.5}{\metre}$, about \SI{26}{\metre}; $\int_0^2 1.5t\,\dd t + \int_2^8 3\,\dd t + \int_8^{11} (3 - (t - 8))\,\dd t = 3 + 18 + 4.5$. (c) A parabola bending upward from a horizontal start, then a straight line of slope \SI{3.0}{\metre\per\second}, then a parabola bending downward and flattening to horizontal at \SI{25.5}{\metre}.
  \item[2.53] (a) $t = \SI{3}{\second}$, at the top of the arch, $x = \SI{9}{\metre}$. (b) Negative: the arch bends downward. Positions at whole seconds $0, 5, 8, 9, 8, 5$~m; first differences $5, 3, 1, -1, -3$; second differences all $-2$, so $a_x = \SI{-2}{\metre\per\second\squared}$. \calctag\ $v_x = 6 - 2t$, $a_x = -2$. (c) Displacement $x(5) - x(0) = \SI{5}{\metre}$; distance $9 + 4 = \SI{13}{\metre}$. (d) $v_x$: a straight line from $+6$ at $t = 0$ through zero at $t = \SI{3}{\second}$ to $-4$ at $t = \SI{5}{\second}$; $a_x$: a horizontal line at $-2$.
  \item[2.54] (a) $v_x(3) = (2.0)(3.0) = \SI{6.0}{\metre\per\second}$; $v_x(9) = 6.0 + (-1.0)(6.0) = 0$. The graph rises in a straight line to $(3, 6)$ then falls in a straight line to $(9, 0)$. (b) At $t = \SI{9.0}{\second}$: the velocity is positive until then and zero there, so the object never turns back. (c) Triangle of base \SI{9.0}{\second} and height \SI{6.0}{\metre\per\second}: \SI{27}{\metre}. \calctag\ $\int_0^3 2t\,\dd t + \int_3^9 (9 - t)\,\dd t = 9 + [9t - \tfrac12 t^2]_3^9 = 9 + (40.5 - 22.5) = \SI{27}{\metre}$. (d) A parabola bending upward from rest to $(3, 9)$, joined smoothly (slope \SI{6}{\metre\per\second} on both sides) to a parabola bending downward that flattens to horizontal at $(9, 27)$; the join is a corner on the velocity--time graph but not on the position--time graph.
  \item[2.55] (a) True: constant slope means constant velocity. (b) True: the velocity changes sign there. (c) False: position zero says nothing about velocity; the thrown ball of Figure~\ref{fig:thrown} returns to $y = 0$ at \SI{15}{\metre\per\second}. (d) False: it is the change in \emph{velocity}; displacement is the area under the velocity--time graph. (e) True: a changing slope means a changing velocity. (f) False: the graph is height against time and is an arch; the path is a vertical line, up and back down the same way.
